% Partie : théorie de la démonstration
% Chapitre : recherche de preuves
% Section : tableaux

\documentclass[../../../include/open-logic-section]{subfiles}

\begin{document}

\olfileid[fr]{pt}{ps}{tab}

\olsection{Tableaux}

Les systèmes de tableaux sont des systèmes de preuve
étroitement liés aux calculs des séquents. Ils conviennent
particulièrement bien à la recherche de preuves, à la
main comme dans un programme. Un tableau est un arbre
de formules plutôt qu'un arbre de séquents. Ses règles
reflètent celles du calcul des séquents, à la différence
des règles de déduction naturelle. Construire un tableau
revient à chercher explicitement un modèle ; on parle
donc aussi de \emph{tableaux sémantiques} ou
d'\emph{arbres de vérité}. Une preuve est obtenue
lorsque cette recherche clôt toutes les branches.

Un tableau \emph{signé} est un arbre de formules
signées, de la forme $\sFmla{\True}{!A}$ ou
$\sFmla{\False}{!A}$. L'arbre se développe vers
le bas. Les formules initiales expriment ce que
l'on cherche à prouver. Pour le séquent
$!A, !B \Sequent !C, !D$, on commence par
$\sFmla{\True}{!A}, \sFmla{\True}{!B},
\sFmla{\False}{!C}, \sFmla{\False}{!D}$.
Une structure dans laquelle $!A$ et $!B$ sont
vraies et $!C$ et $!D$ fausses constitue un
contre-modèle de $!A, !B \Sequent !C, !D$.

On prolonge une feuille, située au bas de l'arbre,
par des règles qui ajoutent des nœuds sur la même
branche ou créent deux branches distinctes.
La formule signée à laquelle on applique la
règle peut se trouver n'importe où au-dessus
de la feuille, sur cette même branche.
Les règles du système classique \Log{Tc}
figurent dans le \olref{tab:Tc}.

\subfile{rules-Tc}

Ces règles correspondent aux règles de \Log{G3c}
lues à rebours. Les signes $\True$ et $\False$
indiquent respectivement une formule principale
dans l'antécédent ou dans le conséquent.
Dans le tableau, les formules antérieures restent
visibles sur la branche après leur traitement.

Une branche est \emph{close} si elle contient
à la fois $\sFmla{\True}{!A}$ et
$\sFmla{\False}{!A}$, ou si elle contient
l'absurde marqué vrai. Un tableau est clos
lorsque toutes ses branches le sont.
Voici un tableau clos pour
$!A \lor (!B \land !C) \Sequent (!A \lor !B) \land
(!A \lor !C)$, commençant par
$\sFmla{\True}{!A \lor (!B \land !C)}$
et $\sFmla{\False}{(!A \lor !B) \land (!A \lor !C)}$ :
\begin{oltableau}
  [\sFmla{\True}{\formula{A} \lor (\formula{B} \land \formula{C})},
    just=\TAss, checked
    [\sFmla{\False}{(\formula{A} \lor \formula{B}) \land
        (\formula{A} \lor \formula{C})}, just=\TAss, checked
      [\sFmla{\True}{\formula{A}}, just={\TRule{\True}{\lor}[1]}
        [\sFmla{\False}{\formula{A} \lor \formula{B}},
          just={\TRule{\False}{\land}[2]}, checked
          [\sFmla{\False}{\formula{A}},
            just={\TRule{\False}{\lor}[4]}
            [\sFmla{\False}{\formula{B}},
              just={\TRule{\False}{\lor}[4]}, close]
          ]
        ]
        [\sFmla{\False}{\formula{A} \lor \formula{C}},
          just={\TRule{\False}{\land}[2]}, checked
          [\sFmla{\False}{\formula{A}},
            just={\TRule{\False}{\lor}[4]}
            [\sFmla{\False}{\formula{C}},
              just={\TRule{\False}{\lor}[4]}, close]
          ]
        ]
      ]
      [\sFmla{\True}{\formula{B} \land \formula{C}},
        just={\TRule{\True}{\lor}[1]}, checked
        [\sFmla{\False}{\formula{A} \lor \formula{B}},
          just={\TRule{\False}{\land}[2]}, checked
          [\sFmla{\True}{\formula{B}},
            just={\TRule{\True}{\land}[3]}, move by=2
            [\sFmla{\True}{\formula{C}},
              just={\TRule{\True}{\land}[3]}
              [\sFmla{\False}{\formula{A}},
                just={\TRule{\False}{\lor}[4]}
                [\sFmla{\False}{\formula{B}},
                  just={\TRule{\False}{\lor}[4]},close
                ]
              ]
            ]
          ]
        ]
        [\sFmla{\False}{\formula{A} \lor \formula{C}},
          just={\TRule{\False}{\land}[2]}, checked
          [\sFmla{\True}{\formula{B}},
            just={\TRule{\True}{\land}[3]}, move by=2
              [\sFmla{\True}{\formula{C}},
                just={\TRule{\True}{\land}[3]}
                [\sFmla{\False}{\formula{A}},
                  just={\TRule{\False}{\lor}[4]}
                  [\sFmla{\False}{\formula{C}},
                  just={\TRule{\False}{\lor}[4]},close
                ]
              ]
            ]
          ]
        ]
      ]
    ]
  ]
\end{oltableau}

Une coche indique que la règle correspondante
a été appliquée à la formule sur toutes les
branches concernées situées au-dessous.
Le symbole $\otimes$ indique une branche close.
Les numéros après les règles désignent la ligne
de la formule signée à laquelle la règle
est appliquée, sur la même branche.

À l'étape $0$, on écrit les formules initiales en haut
du tableau. Sans formule initiale, on part d'une branche
vide. À l'étape $n+1$, pour chaque branche encore
ouverte, on cherche, de haut en bas, la première
formule signée non atomique encore à traiter,
hors des deux formes quantifiées réutilisables.
Si elle existe, on applique sa règle. On coche
une formule ordinaire après l'avoir traitée
sur toutes les branches ouvertes concernées.
Le tableau ci-dessus illustre ces développements
propositionnels.

Les formes $\sFmla{\True}{\lforall}$ et
$\sFmla{\False}{\lexists}$ ne sont pas cochées
définitivement : elles demandent un traitement
particulier. Les règles $\TRule{\True}{\forall}$
et $\TRule{\False}{\lexists}$ doivent produire
toutes les instances requises par les termes
clos. À chaque étape, après le traitement
ordinaire, on applique donc la règle, sur chaque
branche restée ouverte, aux occurrences réutilisables
déjà présentes au début de l'étape. On conserve pour
chacune, sur cette branche, l'historique des termes
déjà employés.

Le terme $t$ choisi est le premier terme~$t$
non encore employé pour cette occurrence,
construit à partir des constantes de la branche
et des symboles de fonction des formules
initiales. L'ensemble associé contient au
départ les constantes initiales, ou $\Obj c_0$
s'il n'y en a aucune ; il conserve ensuite
les constantes ajoutées. On ajoute l'instance
$\sFmla{\True}{!B(t)}$ ou~$\sFmla{\False}{!B(t)}$,
même si elle apparaît déjà sur la branche.
Sans fonction d'arité positive, on teste si la liste
finie des constantes a été épuisée. Dans ce cas, on
passe à l'occurrence suivante. Avec une telle fonction,
le parcours trouve toujours un terme encore inutilisé.
Les règles à constante propre choisissent
une constante absente de l'ensemble associé,
et l'y ajoutent. Les historiques et ensembles
sont copiés lors des ramifications.

La recherche s'arrête avec succès dès que toutes
les branches sont closes. Si une étape n'ajoute
aucun nœud et qu'une branche reste ouverte,
elle s'arrête sans succès : une telle branche
est finie et saturée.

Si cette recherche ne produit pas un tableau
clos, elle fournit soit une branche ouverte
finie saturée, soit un arbre infini à
ramifications finies, qui possède une branche
infinie. Dans le cas fini, toutes les règles
ordinaires ont été appliquées et toutes les
instances réutilisables sur le domaine de
termes ont été produites. Dans le cas infini,
l'équité et l'énumération assurent la même
saturation. Comme dans la \olref[cpl]{sec},
on définit une structure~$\Struct{M}$ telle que
$\Sat{M}{!A}$ lorsque $\sFmla{\True}{!A}$
figure sur la branche, et $\Sat/{M}{!A}$
lorsque $\sFmla{\False}{!A}$ y figure.
On prend
$\Theta = \Setabs{!A}{\sFmla{\True}{!A} \text{ sur la branche}}$
et
$\Xi = \Setabs{!A}{\sFmla{\False}{!A} \text{ sur la branche}}$.

Pour traduire un tableau clos en une preuve
de \Log{G3c}, on associe des séquents aux
blocs de prolongement. Si les formules initiales
correspondent à $\Gamma \Sequent \Delta$,
chacun des nœuds initiaux reçoit le même
séquent $\Gamma \Sequent \Delta$.
Lorsqu'on prolonge la branche sur
$\sFmla{\True}{!A}$, son séquent a la forme
$!A, \Pi \Sequent \Lambda$ ; pour
$\sFmla{\False}{!A}$, il a la forme
$\Pi \Sequent \Lambda, !A$.
Une règle \TRule{\True}{\star} correspond
à \LeftR{\star}, avec la même formule
principale ; une règle \TRule{\False}{\star}
correspond à \RightR{\star}.
Les prémisses sont les séquents associés
aux branches après l'ajout du bloc de formules.
À une clôture sur $\sFmla{\True}{!A}$
et $\sFmla{\False}{!A}$ correspond un
séquent $!A, \Pi \Sequent \Lambda, !A$.
Plusieurs nœuds successifs d'un même bloc
peuvent recevoir le même séquent ; on
supprime ces répétitions dans la preuve.

La clôture sur $\sFmla{\True}{\lfalse}$
%
correspond à l'axiome de l'absurde de \Log{G3c}.
Une clôture sur une formule composée utilise
la preuve d'identité de la
\olref[pt][seq][ex]{prop:G3c-axioms}.
Ces deux conventions sont nécessaires pour
que « tableau clos » donne une preuve du
calcul avec ses seuls axiomes atomiques.

Pour une conversion générale, on garde
dans le séquent de branche toutes les
formules déjà présentes, y compris celles
qui ont été traitées, avec leurs multiplicités.
Une expansion ordinaire ajoute ses formules
actives tout en laissant la principale dans
le contexte. En appliquant la règle logique
de \Log{G3c} aux preuves des branches ainsi
prolongées, on obtient donc deux copies
de la principale dans la conclusion.
On les contracte par la
\olref[pt][seq][inv]{prop:G3c-cont-adm}.
Pour l'universelle vraie et l'existentielle
fausse, la règle de \Log{G3c} conserve déjà
la principale ; aucune duplication
supplémentaire n'est nécessaire.
Les règles à constante propre satisfont
leur condition, puisque la constante choisie
est absente de toute la branche antérieure.

Une récurrence, depuis les feuilles closes
vers la racine du tableau fini, fournit ainsi
les preuves de tous les séquents de branche.
La contraction est admissible : on remplace
chacune de ses utilisations par la transformation
de preuves correspondante. Le résultat final
emploie uniquement les règles de \Log{G3c}.
Les contractions réconcilient la persistance des
formules dans le tableau avec leur traitement
%
dans le calcul des séquents.

Pour l'exemple précédent, on peut obtenir
la preuve plus concise suivante. Si les
formules aux feuilles d'identité ne sont
pas atomiques, ces feuilles sont à développer
en leurs preuves d'identité.
{\small\[
\Axiom$!A \fCenter !A, !B$
\RightLabel{\RightR{\lor}}
\UnaryInf$!A \fCenter !A \lor !B$
\Axiom$!A \fCenter !A, !C$
\RightLabel{\RightR{\lor}}
\UnaryInf$!A \fCenter !A \lor !C$
\RightLabel{\RightR{\land}}
\BinaryInf$!A \fCenter (!A \lor !B) \land (!A \lor !C)$
\Axiom$!B, !C \fCenter !A, !B$
\RightLabel{\RightR{\lor}}
\UnaryInf$!B, !C \fCenter !A \lor !B$
\RightLabel{\LeftR{\land}}
\UnaryInf$!B \land !C \fCenter !A \lor !B$
\Axiom$!B, !C \fCenter !A, !C$
\RightLabel{\RightR{\lor}}
\UnaryInf$!B, !C \fCenter !A \lor !C$
\RightLabel{\LeftR{\land}}
\UnaryInf$!B \land !C \fCenter !A \lor !C$
\RightLabel{\RightR{\land}}
\BinaryInf$!B \land !C \fCenter (!A \lor !B) \land
(!A \lor !C)$
\RightLabel{\LeftR{\lor}}
\BinaryInf$!A \lor (!B \land !C) \fCenter (!A \lor !B) \land
(!A \lor !C)$
\DisplayProof
\]}

\end{document}
